\section{Method}
\label{sec:method}

We describe PALETTE (\S\ref{subsec:m-transform}), how it is applied during
training (\S\ref{subsec:m-training}), and the causal ablation used to
isolate texture preservation as the source of its effect
(\S\ref{subsec:m-ablation}).

\subsection{The PALETTE transform}
\label{subsec:m-transform}

Let $x$ denote the intensity volume of a training scan, min-max normalised
to $[0,1]$, and $F \subset \Omega$ its foreground voxels (an intensity
threshold; $x_v > 0.01$ in our implementation). PALETTE applies four steps
at every training iteration, illustrated end-to-end in
\cref{fig:method-pipeline} on one real CHAOS (abdominal MRI) and one real
OnHarmony (brain MRI) volume.

\paragraph{1. Intensity partition.} We run 1-D $k$-means on the foreground
intensities $\{x_v : v \in F\}$ with $k$ sampled per iteration from
$\{2,\dots,6\}$, producing $k$ intensity classes $C_1, \dots, C_k$
(\cref{fig:method-pipeline}, column 2). This partition is unsupervised and
requires no anatomical label map: it groups voxels by intensity alone, so
the same operator applies unmodified to CT, T1, T2, FLAIR, or any other
contrast. With probability $0.10$ this step is skipped and a single global
remap is applied instead (step 3 below, with $C=1$).

\paragraph{2. Spatial sub-parcellation.} Each intensity class $C_i$ is
further subdivided into spatial sub-regions by a Voronoi tessellation
seeded with points sampled within $C_i$ (\cref{fig:method-pipeline}, column
3 shows the combined effect of this step and the remap of step 3), giving
sub-regions $R_{i,1}, \dots, R_{i,m_i}$ that are contiguous in space as
well as similar in intensity. This second partition breaks the spatial
correlation of a purely intensity-based grouping: two voxels with the same
intensity in different parts of the anatomy can land in different
sub-regions and receive independent contrast perturbations. Each class is
subdivided into $m_i$ sub-regions sampled from $\{2,\dots,10\}$ with
probability $0.60$, and left as one region (no sub-split) with probability
$0.40$.

\paragraph{3. Signed affine remap.} For each sub-region $R$, let $\bar{x}_R$
denote its mean intensity, $\mu_R \sim U(0,1)$ a resampled target mean, and
draw $\alpha_R \sim \pm U(0.5, 2.0)$ (a magnitude in $[0.5, 2.0]$ with an
independently sampled sign). Every voxel $v \in R$ is remapped as
\begin{equation}
  y_v = \mu_R + \alpha_R \, (x_v - \bar{x}_R).
  \label{eq:affine}
\end{equation}
Because \cref{eq:affine} is a per-region affine function of the
\emph{real} intensity $x_v$, it preserves the relative structure of the
signal within each region: edges, gradients, and partial-volume transitions
between adjacent voxels are scaled, not replaced. Only the region-to-region
contrast relationship is randomized, and $\alpha_R < 0$ additionally allows
local contrast inversions that no acquisition protocol produces on its own,
widening the space of contrasts the network sees beyond what any single
real sequence could supply. Background voxels are zeroed at this step (and
remain zero through the rest of the pipeline). The result is optionally
Gaussian-blurred (a $\sigma$ sampled from $\{0,0,0,0.3,0.5,0.8\}$, weighted
toward no blur).

\paragraph{4. Anatomical-label affine remap (optional).} Steps 1--3 are
unsupervised: the $k$-means/Voronoi partition groups voxels by intensity
and spatial proximity, which need not coincide with the true anatomical
boundary between two structures. When two adjacent tissues happen to
share similar intensity at a given training iteration, the partition can
place both sides of a real label border inside the same region, so step 3
assigns them the same target contrast and the boundary is, by chance,
erased rather than randomized. This last step corrects for that: for each
real anatomical label $\ell$ present in the volume, with probability
$0.5$ (skipped below $4$ foreground voxels) we resample a fresh
$(\mu_\ell, \alpha_\ell)$ and apply \cref{eq:affine} \emph{again},
restricted to $\ell$'s voxels only, \textbf{on top of the step-3 output}
-- i.e.\ this affine remap acts on the already-remapped (step 1--3)
intensities, not on the original voxel values, so it composes with rather
than replaces the unsupervised remap (\cref{fig:method-pipeline}, column
4 shows the resulting per-label flat-remap value, computed prior to this
step). Because the label assignment is independent of whatever
$k$-means/Voronoi region a voxel fell into, this step reintroduces a
full-strength contrast discontinuity exactly at the true label border
regardless of whether the unsupervised partition happened to preserve it
-- sharpening the boundary cue the segmentation network can rely on,
rather than sharpening the underlying anatomy. The step only ever
consumes the segmentation labels the network is already being trained to
predict -- whatever labels happen to be annotated in the training set,
including sparse or partial annotations -- so it requires no additional
segmentation effort beyond what training already assumes. It is also
optional: it can be disabled entirely (\texttt{label\_remap\_prob}$=0$),
in which case PALETTE degrades gracefully to steps 1--3 alone, e.g.\ on
an unlabelled dataset. A final foreground z-score normalisation is
applied after an optional second blur pass (same $\sigma$ distribution as
step 3), giving the training input shown in the last column of
\cref{fig:method-pipeline}.

\smallskip\noindent\emph{Per-step stochasticity, summarised:} step 1 is
replaced by a single global remap with probability $0.10$; each
intensity class is Voronoi-split with probability $0.60$ (else left as
one region); step 3 is always applied to whatever partition results from
steps 1--2; step 4 is applied independently per label with probability
$0.5$.

\begin{figure*}[t]
  \centering
  \includegraphics[width=\linewidth]{figures/palette_method_figure}
  \caption{The PALETTE pipeline applied to one real CHAOS T2SPIR volume
  (top) and one real OnHarmony T1w volume (bottom). For display only, every
  panel is min--max normalised to $[0,1]$ (the foreground z-score used at
  training time is never applied for visualisation); no panel is generated
  by an external model -- every value shown is a deterministic function of
  the same real volume. Column 1: input. Column 2: \emph{$k$-means flat
  remap} -- each $k$-means class is painted with the precomputed target
  mean $\mu$ of \cref{eq:affine} for its dominant (largest) Voronoi
  sub-region, i.e.\ the flat value that sub-region would take under the
  transform, not an average of observed pixel intensities. Column 3: the
  same flat-remap logic applied after Voronoi sub-parcellation, so each
  sub-region is painted with its own precomputed $\mu$. Column 4: each
  real anatomical label is additionally painted with its own precomputed
  $\mu$, overwriting the column-3 value on that label's footprint. Column
  5: final PALETTE training input, obtained by applying the signed affine
  remap of \cref{eq:affine} with texture (not a flat fill) at both the
  region level and, where labels are available, the anatomical-label
  level. Columns 2--4 isolate the flat target value chosen \emph{before}
  the transform is applied with texture; column 5 shows the texture the
  transform actually produces once that value is used as its mean.}
  \label{fig:method-pipeline}
\end{figure*}

\subsection{Training protocol}
\label{subsec:m-training}

PALETTE is applied on-the-fly as one contrast-augmentation operator inside
an existing GPU-accelerated nnU-Net augmentation
pipeline~\cite{molinier2026auglab,isensee2021nnunet}, alongside that
pipeline's standard spatial and artifact augmentations; it adds no learned
parameters and no measurable training-time overhead. A single
configuration -- the same $k$-range, sub-region count, blur schedule, and
$\alpha$ range given in \S\ref{subsec:m-transform} -- is tuned once and
reused unchanged across every dataset and modality in
\cref{sec:experiments}; no dataset receives its own hyperparameter sweep.

\subsection{Causal ablation: isolating texture preservation}
\label{subsec:m-ablation}

PALETTE differs from label-generative synthesis (e.g.\ SynthSeg) along two
axes at once: the partition it perturbs (unsupervised $k$-means $+$ Voronoi
vs.\ a supplied anatomical label map as the sole generative substrate) and
the perturbation itself (an affine remap of real intensities vs.\
resampling from a per-label Gaussian mixture). Comparing PALETTE to
SynthSeg directly therefore cannot attribute an advantage to either axis in
isolation.

We construct a controlled variant, \emph{noise-fill}, that holds the
$k$-means/Voronoi partition and the per-region statistics $(\mu_R,
\alpha_R)$ from \cref{eq:affine} exactly fixed, but fills each sub-region
with samples drawn from a Gaussian with that region's $(\mu_R, \alpha_R
\sigma_R)$ instead of remapping the real intensities inside it -- i.e.\ the
same partitioning scheme, with SynthSeg-style noise substituted for the
affine remap. Because the partition and target statistics are identical to
PALETTE's own, any difference in downstream performance between PALETTE and
noise-fill isolates the one variable that changed: whether real texture is
preserved or discarded within each region. We report this comparison
separately for texture-defined structures (identified by fine internal
gradients or low boundary contrast) and boundary-defined structures
(identified by a sharp, high-contrast border) in \cref{sec:experiments}, as
the texture-preservation hypothesis predicts a widening gap specifically on
the former.
